\documentclass{exam}
\usepackage{../../shah350}

\hypersetup{colorlinks=true, linktoc=section, linkcolor=blue}

\pagestyle{headandfoot}
\firstpageheadrule
\runningheadrule
\firstpageheader{Prof. Lepowsky \\ Linear Algebra}{Homework \#6}{\href{https://github.com/jeevanshah07/MATH350/blob/main/homework/homework6/main.tex}{\underline{Source}} \\ Jeevan Shah}
\runningheader{Linear Algebra \\ Homework \#6}{}{Shah}
\firstpagefooter{}{\thepage}{}
\runningfooter{ }{\thepage}{ }

\printanswers


\begin{document}
  \colorbox{red}{\underline{$2.4.15$}:}
  \begin{solution}
    \begin{proof}
      $(\Rightarrow)$ Suppose that $T$ is an isomorphism. Since $T$ is injective, $T(\beta)$ contains exactly $n$ vectors, so it is sufficient to prove that $T(\beta)$ is linearly independent. Write $\beta = \braces{b_1, \dots, b_n}$ and suppose $\sum_{i=1}^{n} a_i T(b_i) = \vec{0}$. By linearity, $T\left(\sum_{i=1}^{n} a_i b_i\right) = \vec{0}$, so by injectivity $\sum_{i=1}^{n} a_i b_i = \vec{0}$. Since $\beta$ is linearly independent, every $a_i = 0$. Thus, $T(\beta)$ is linearly independent and thus a basis for $W$.
      $(\Leftarrow)$ Suppose $T(\beta)$ is a basis for $W$. It is sufficient to show that $T$ is a bijection to prove that it is an isomorphism, as if $T$ is bijective, then it is invertible. It is clear that $T$ is onto: since $T(\beta)$ spans $W$, every $w \in W$ can be written $w = \sum a_i T(b_i) = T\left(\sum a_i b_i\right)$ and so lies in the image of $T$. Now, since $\dim V = \dim W$, by Thm~2.5, $T$ is also one-to-one. Thus $T$ is an isomorphism. \\
    \end{proof}
  \end{solution}

  \colorbox{red}{\underline{$2.5.2$}:}
  \begin{solution}
    \begin{parts}
      \setcounter{partno}{1}
      \part We have that 
      \[
        \begin{pmatrix}
          0 \\ 10
        \end{pmatrix}
        = 4\begin{pmatrix}
          -1 \\ 3
        \end{pmatrix}
        + 2\begin{pmatrix}
          2 \\ -1
        \end{pmatrix}
        \quad\text{and}\quad
        \begin{pmatrix}
          5 \\ 0 
        \end{pmatrix}
        = \begin{pmatrix}
          -1 \\ 3
        \end{pmatrix}
        + 3\begin{pmatrix}
          2 \\ -1
        \end{pmatrix}
      \]
      so, 
      \[
        \boxed{
          Q = \begin{pmatrix}
            4 & 1 \\
            2 & 3
          \end{pmatrix}
        }
      \]

      \part We have that 
      \[
          \begin{pmatrix}
            1 \\ 0
          \end{pmatrix}
          = 3\begin{pmatrix}
            2 \\ 5
          \end{pmatrix}
          + 5\begin{pmatrix}
            -1 \\ -3
          \end{pmatrix}
          \quad\text{and}\quad
          \begin{pmatrix}
            0 \\ 1
          \end{pmatrix}
          = -\begin{pmatrix}
            2 \\ 5
          \end{pmatrix}
          -2 \begin{pmatrix}
            -1 \\ -3
          \end{pmatrix}
      \]
      so
      \[
          \boxed{
            Q = \begin{pmatrix}
              3 & -1 \\
              5 & -2
            \end{pmatrix}
          }
      \]
    \end{parts}
  \end{solution}

  \colorbox{red}{\underline{$2.5.3c$}:}
  \begin{solution}
    Since $P_2(\RR) \cong \RR^3$ we know 
    \[
      \beta = \braces{
        \begin{pmatrix}
          0 \\ -1 \\ 2
        \end{pmatrix},
        \begin{pmatrix}
          1 \\ 0 \\ 3
        \end{pmatrix},
        \begin{pmatrix}
          0 \\ 0 \\ 1
        \end{pmatrix}
      }
      \quad\text{and}\quad
      \beta' = \braces{
        \begin{pmatrix}
          1 \\ 0 \\ 0
        \end{pmatrix},
        \begin{pmatrix}
          0 \\ 1 \\ 0
        \end{pmatrix},
        \begin{pmatrix}
          0 \\ 0 \\ 1
        \end{pmatrix}
      }
    \]
    and, 
    \[
        \begin{pmatrix}
          1 \\ 0 \\ 0
        \end{pmatrix} 
        = 0 \begin{pmatrix}
          0 \\ -1 \\ 2
        \end{pmatrix}
        + \begin{pmatrix}
          1 \\ 0 \\ 3
        \end{pmatrix}
        - 3 \begin{pmatrix}
          0 \\ 0 \\ 1
        \end{pmatrix}
        \quad\text{and}\quad
        \begin{pmatrix}
          0 \\ 1 \\ 0
        \end{pmatrix}
        = - \begin{pmatrix}
          0 \\ -1 \\ 2
        \end{pmatrix}
        + 0\begin{pmatrix}
          1 \\ 0 \\ 3
        \end{pmatrix}
        + 2 \begin{pmatrix}
          0 \\ 0 \\ 1
        \end{pmatrix}
    \]
    so, 
    \[
      \boxed{
        Q = \begin{pmatrix}
          0 & -1 & 0 \\
          1 & 0 & 0 \\
          -3 & 2 & 1
        \end{pmatrix}
      }
    \]
  \end{solution}

  \colorbox{red}{\underline{$2.5.5$}:}
  \begin{solution}
    Directly applying the theorem we see 
    \[
      [T]_{\beta'} = \begin{pmatrix}
        1/2 & 1/2 \\
        1/2 & -1/2
      \end{pmatrix}
      \begin{pmatrix}
        0 & 1 \\
        0 & 0
      \end{pmatrix}
      \begin{pmatrix}
        1 & 1 \\
        1 & -1 
      \end{pmatrix}
      = \boxed{
        \begin{pmatrix}
          1/2 & -1/2 \\
          1/2 & -1/2
        \end{pmatrix}
      }
    \]
  \end{solution}

  \colorbox{red}{\underline{$2.5.6$}:}
  \begin{solution}
    \begin{parts}
      \setcounter{partno}{1}
      \part We first find $Q$ by letting $Q$ be the $2 \times 2$ matrix whose $j$th column is the $j$th vector of $\beta$:
      \[
        \boxed{
            Q = \begin{pmatrix}
            1 & 1 \\
            1 & -1
          \end{pmatrix}
        }
      \]
      Then, by the corollary to Thm~2.23:
      \[
          [L_A]_\beta = Q^{-1}AQ = \frac{1}{2}\begin{pmatrix}
            1 & 1 \\
            1 & -1 
          \end{pmatrix}
          \begin{pmatrix}
            1 & 2 \\
            2 & 1 
          \end{pmatrix}
          \begin{pmatrix}
            1 & 1 \\
            1 & -1 
          \end{pmatrix}
          = \boxed{
            \begin{pmatrix}
              3 & 0 \\
              0 & -1
            \end{pmatrix}
          }
      \]

      \part As above we have 
      \[
        \boxed{
          Q = \begin{pmatrix}
            1 & 1 & 1 \\
            1 & 0 & 1 \\
            1 & 1 & 2
          \end{pmatrix}
        }
      \]
      so, 
      \begin{align*}
        (Q \mid I_3) = \begin{Lamatrix}{6}{3}
          1 & 1 & 1 & 1 & 0 & 0 \\
          1 & 0 & 1 & 0 & 1 & 0 \\
          1 & 1 & 2 & 0 & 0 & 1
        \end{Lamatrix}
        \xrightarrow{r_2 \leftrightarrow r_1}&
        \begin{Lamatrix}{6}{3}
          1 & 0 & 1 & 0 & 1 & 0 \\
          1 & 1 & 1 & 1 & 0 & 0 \\
          1 & 1 & 2 & 0 & 0 & 1
        \end{Lamatrix} \\
        \xrightarrow{r_3 \to r_3 - r_2}& \begin{Lamatrix}{6}{3}
          1 & 0 & 1 & 0 & 1 & 0 \\
          1 & 1 & 1 & 1 & 0 & 0 \\
          0 & 0 & 1 & -1 & 0 & 1
        \end{Lamatrix} \\
        \xrightarrow{r_2 \to r_2 - r_1}& \begin{Lamatrix}{6}{3}
          1 & 0 & 1 & 0 & 1 & 0 \\
          0 & 1 & 0 & 1 & -1 & 0 \\
          0 & 0 & 1 & -1 & 0 & 1
        \end{Lamatrix} \\
        \xrightarrow{r_1 \to r_1 - r_3}& \begin{Lamatrix}{6}{3}
          1 & 0 & 0 & 1 & 1 & -1 \\
          0 & 1 & 0 & 1 & -1 & 0 \\
          0 & 0 & 1 & -1 & 0 & 1
        \end{Lamatrix}
        = (I_3 \mid Q^{-1})
      \end{align*}
      so 
      \[
        [L_A]_\beta = \begin{pmatrix}
          1 & 1 & -1 \\
          1 & -1 & 0 \\
          -1 & 0 & 1
        \end{pmatrix}
        \begin{pmatrix}
          1 & 1 & -1 \\
          2 & 0 & 1 \\
          1 & 1 & 0
        \end{pmatrix}
        \begin{pmatrix}
          1 & 1 & 1 \\
          1 & 0 & 1 \\
          1 & 1 & 2
        \end{pmatrix}
        = \boxed{
          \begin{pmatrix}
            2 & 2 & 2 \\
            -2 & -3 & -4 \\
            1 & 1 & 2
          \end{pmatrix}
        }
      \]
    \end{parts}
  \end{solution}

  \newpage 

  \colorbox{red}{\underline{$3.2.5a$}:}
  \begin{solution}
    We have 
    \begin{align*}
      \begin{pmatrix}
        1 & 2 & 1 \\
        -1 & 1 & 2 \\
        1 & 0 & 1
      \end{pmatrix}
      \xrightarrow[r_3 \to r_3 - r_1]{r_2 \to r_2 + r_1}& \begin{pmatrix}
        1 & 2 & 1 \\
        0 & 3 & 3 \\
        0 & -2 & 0 
      \end{pmatrix} \\
      \xrightarrow[r_1 \to r_1 - 2r_2]{r_2 \to \frac{1}{3}r_2}& \begin{pmatrix}
        1 & 0 & -1 \\
        0 & 1 & 1 \\
        0 & -2 & 0
      \end{pmatrix} \\
      \xrightarrow[r_2 \to -\frac{1}{2}r_2]{r_2 \leftrightarrow r_3}& \begin{pmatrix}
        1 & 0 & -1 \\
        0 & 1 & 0 \\
        0 & 1 & 1 
      \end{pmatrix} \\
      \xrightarrow{r_3 \to r_3 - r_2}& \begin{pmatrix}
        1 & 0 & -1 \\
        0 & 1 & 0 \\
        0 & 0 & 1
      \end{pmatrix} \\
      \xrightarrow{r_1 \to r_1 + r_3}& \begin{pmatrix}
        1 & 0 & 0 \\
        0 & 1 & 0 \\
        0 & 0 & 1
      \end{pmatrix}
    \end{align*}
    so, $\rank(A) = 3$. Thus, 
    \begin{align*}
      (A \mid I_3) = \begin{Lamatrix}{6}{3}
        1 & 2 & 1 & 1 & 0 & 0 \\
        -1 & 1 & 2 & 0 & 1 & 0 \\
        1 & 0 & 1 & 0 & 0 & 1
      \end{Lamatrix}
      \xrightarrow{}
      \begin{Lamatrix}{6}{3}
        1 & 0 & 0 & 1/6 & -1/3 & 1/2 \\
        0 & 1 & 0 & 1/2 & 0 & -1/2 \\
        0 & 0 & 1 & -1/6 & 1/3 & 1/2
      \end{Lamatrix} = (I_3 \mid A^{-1})
    \end{align*}
    by the same sequence of elementary row operations as above. Hence, 
    \[
      \boxed{
        A^{-1} = \begin{pmatrix}
          1/6 & -1/3 & 1/2 \\
          1/2 & 0 & -1/2 \\
          -1/6 & 1/3 & 1/2
        \end{pmatrix}
      }
    \]
  \end{solution}

  \colorbox{red}{\underline{$3.2.6a$}:}
  \begin{solution}
    Let $\beta$ be the standard ordered basis of $P_2(\RR)$ so 
    \[
      [T]_\beta = \begin{pmatrix}
        -1 & 2 & 2 \\
        0 & -1 & 4 \\
        0 & 0 & -1
      \end{pmatrix}
    \]
    since $\det([T]_\beta) \neq 0$, it follows that $([T]_\beta)^{-1}$ exists.
    \begin{align*}
      ([T]_\beta \mid I_3) = \begin{Lamatrix}{6}{3}
        -1 & 2 & 2 & 1 & 0 & 0 \\
        0 & -1 & 4 & 0 & 1 & 0 \\
        0 & 0 & -1 & 0 & 0 & 1
      \end{Lamatrix}
      \xrightarrow[r_3 \to -r_3]{r_1 \to -r_1,\ r_2 \to -r_2}&
      \begin{Lamatrix}{6}{3}
        1 & -2 & -2 & -1 & 0 & 0 \\
        0 & 1 & -4 & 0 & -1 & 0 \\
        0 & 0 & 1 & 0 & 0 & -1
      \end{Lamatrix} \\
      \xrightarrow[r_1 \to r_1 + 2r_3]{r_2 \to r_2 + 4r_3}& \begin{Lamatrix}{6}{3}
        1 & -2 & 0 & -1 & 0 & -2 \\
        0 & 1 & 0 & 0 & -1 & -4 \\
        0 & 0 & 1 & 0 & 0 & -1
      \end{Lamatrix} \\
      \xrightarrow{r_1 \to r_1 + 2r_2}& \begin{Lamatrix}{6}{3}
        1 & 0 & 0 & -1 & -2 & -10 \\
        0 & 1 & 0 & 0 & -1 & -4 \\
        0 & 0 & 1 & 0 & 0 & -1
      \end{Lamatrix}
      = (I_3 \mid ([T]_\beta)^{-1})
    \end{align*}
    so, 
    \[
      \boxed{
        ([T]_\beta)^{-1} = \begin{pmatrix}
          -1 & -2 & -10 \\
          0 & -1 & -4 \\
          0 & 0 & -1
        \end{pmatrix}
      }
    \]
    and 
    \[
      \boxed{
        T^{-1}(a_0 + a_1 x + a_2 x^2) = (-a_0 - 2a_1 - 10a_2) + (-a_1 - 4a_2)x - a_2x^2
      }
    \]
  \end{solution}

  \colorbox{red}{\underline{$3.4.6$}:}
  \begin{solution}
    Notice that the first, third, and fifth columns of $B = \operatorname{rref}(A)$ are $e_1, e_2$, and $e_3$ respectively. Since $b_2 = -3e_1$ it follows that $a_2 = -3a_1$, 
    \[
      a_2 = -3\begin{pmatrix}
        1 \\ -2 \\ -1 \\ 3
      \end{pmatrix} = \begin{pmatrix}
        -3 \\ 6 \\ 3 \\ -9
      \end{pmatrix}
    \]
    Similarly, $b_4 = 4e_1 + 3e_2 \Rightarrow a_4 = 4a_1 + 3a_3$, 
    \[
      a_4 = 4\begin{pmatrix}
        1 \\ -2 \\ -1 \\ 3
      \end{pmatrix}
      + 3\begin{pmatrix}
        -1 \\ 1 \\ 2 \\ -4
      \end{pmatrix}
      = \begin{pmatrix}
        1 \\ -5 \\ 2 \\ 0
      \end{pmatrix}
    \]
    Now, $b_6 = 5e_1 + 2e_2 - e_3 \Rightarrow a_6 = 5a_1 + 2a_3 - a_5 \Rightarrow a_5 = 5a_1 + 2a_3 - a_6$, 
    \[
      a_5 = 5\begin{pmatrix}
        1 \\ -2 \\ -1 \\ 3
      \end{pmatrix}
      + 2\begin{pmatrix}
        -1 \\ 1 \\ 2 \\ -4
      \end{pmatrix}
      -\begin{pmatrix}
        3 \\ -9 \\ 2 \\ 5
      \end{pmatrix}
      = \begin{pmatrix}
        0 \\ 1 \\ -3 \\ 2
      \end{pmatrix}
    \]
    Hence, 
    \[
      \boxed{
        A = \begin{pmatrix}
          1 & -3 & -1 & 1 & 0 & 3 \\
          -2 & 6 & 1 & -5 & 1 & -9 \\
          -1 & 3 & 2 & 2 & -3 & 2 \\
          3 & -9 & -4 & 0 & 2 & 5
        \end{pmatrix}
      }
    \]
  \end{solution}
\end{document}